\section{Machine Learning Background}

\red{As an introduction to the methods used in this field, this section covers the concept of neural posterior approximation, some theory behind the class of models analyzed in this thesis as well as a summary of metric from information theory, and how they can evaluate a model's quality.}

In the following, $x$ denotes a data sample, and $\theta$ a corresponding parameter. Probability densities are denoted $p$, and their approximations with $q$. Functions $s$ denote objects that summarize data, so $s(x)$ is such a summary. Neural networks are denoted with an index $\phi or \psi$, so $q_\phi$ is an estimate of a probability distribution done with a neural network, and $s_\psi$ is a summary done by a neural network.

\subsection{Neural posterior estimation}
In a Bayesian inference task, processing these data would typically be done by choosing:
\begin{enumerate}
    \item   a likelihood $p(x|\theta)$ describing how observed data arised from a specific set of parameters,
    \item   a prior $p(\theta)$ describing beliefs about the parameter distribution before seeing data.
\end{enumerate}
Given data d, one then calculates the posterior distribution using Bayes' theorem:
\begin{align}
    p(\theta|x) \sim p(x|\theta) p(\theta).
\end{align}
In many applied cases, an analytic shape of the likelihood is not available. An example is when the data come from a complex simulator. Then one can sample from this simulator and as such from the likelihood, but the expression $p(x_{\text{observed}}|\theta)$ is not obtainable and as such, the posterior cannot be calculated. In such setting, an alternative is to model the posterior separately, such that given $x_{\text{observed}}$, one can at least sample from the posterior: $\theta \sim q(\theta | x_{\text{observed}})$. This is known as likelihood-free inference. When the posterior is approximated using a neural network $q_\phi(\theta |x)$, one also speaks of neural posterior estimation (NPE).

An advantage is that unlike with the traditional Bayesian methods, the posterior does not have to be re-calculated given new data: the neural network simply has to be evaluated on the new dataset. In this way, the cost of obtaining the posterior was payed beforehand during training. This concept is known as amortized inference. In order to train the model, data is needed, and this is where simulators are implemented. The posterior model is trained on simulated data and once trained, one can sample parameters by $\theta \sim q_\phi(\theta | x_{\text{observed}})$. In this way, NPE becomes a way of implementing simulation-based inference.

\subsection{Information Theory and Neural Estimators}
Traditionally, data as complex as those expected from SKA are analyzed using physics-inspired summary statistics $s_\text{phys}(x)$, which map the data to a lower-dimensional space with known interpretation. An example is the power spectrum. Such a summary comes with the benefit of interpretability, but the caveat is a potentially large information loss. This motivates the use of a neural networks to learn a different summary $s_\phi(x)$, and the procedure is then called representation learning. The summarry is hopefully more informative of the parameter of interest, though potentially less interpretable. If $\mathbb{I}$ is a measure of informativeness between the data and the parameters, this increased informativeness can be expressed as:
\begin{align}
    \mathbb{I}(\theta;x) 
    \geq \mathbb{I}(\theta;s_\phi(x)) 
    > \mathbb{I}(\theta;s_{\text{phys}}(x)).
\end{align}
The necessitates a measure of informativeness for use in representation learning, which is introduced in this section.

\subsubsection*{Some Basic Measures in Information Theory}
Given a random variable $X$ distributed as $p$, $X \sim p$, the entropy of $X$ is defined as:
\begin{align}
\mathbb{H}(X) = -\mathbb{E}_{X} \left[ \log p(X) \right] \geq 0.
\end{align}
If the logarithm is chosen with base 2, one defines the unit of this quantity as bits. If one uses the natural logarithm, the entropy is measured in units of nats. 

Entropy is a measure of the uncertainty of a random variable. A high entropy corresponds to a lack of predictability of the random variable's outcome. As an example, the entropy of the uniform distribution is maximal, while that of the delta distribution is zero. One therefore considers entropy to be a measure of the amount of information contained in the random variable's distribution.

The generalization to conditional entropy is:
\begin{align}
    \mathbb{H}(X|Y) =  -\mathbb{E}_{X, Y} \left[ \log p(X|Y) \right]
\end{align} 

Another important metric in information theory is the Kullback-Leibler divergence. Let $X \sim p$ and $Y \sim q$ be two random variables over the same sample space. The KL-divergence between $p$ and $q$ is defined as:
\begin{align}
\mathbb{D}_{\text{KL}}(p || q) = \mathbb{E}_{X} \left[ \log p(X) - \log q(X) \right] \geq 0
\end{align}
The KL-divergence satisfies the identity axiom, $\mathbb{D}_{\text{KL}}(p || q) = 0 \Leftrightarrow p = q$. It is not symmetric and does not satisfy the triangle inequality and as such is not a metric. However, it does still carry an interpretation as a distance between two probability distributions: from the above expression it is clear that the more $q$ resembles $p$, the smaller the measured KL-divergence becomes. A smaller the KL-divergence between the two distributions means there is less information loss by using  $q$ to approximate $p$.

Since the KL-divergence is assymmetric, one distinguishes $\mathbb{D}_{\text{KL}}(p || q)$ as the forwards KL between $p$ and $q$, and $\mathbb{D}_{\text{KL}}(q || p)$ as the backwards KL between $p$ and $q$. Minimizing the forwards KL-divergence between the data distribution and the model distribution is equivalent to maximizing the negative log-likelihood, because
\begin{align}
    \mathbb{D}_{\text{KL}}(p_{\text{data}} || q_\phi) = 
    -\mathbb{H}(p_{\text{data}}) - \E_{x \sim p_{\text{data}}} 
    \bigl[ \log q_\phi \bigr] = \text{NLL} + \text{const}
\end{align}

Lastly, the mutual information between two random variables $X$ and $Y$ is defined as:
\begin{align}
\mathbb{I}(X; Y)
         &= \mathbb{D}_{\text{KL}}(p(X, Y) || p_X(X) p_Y(Y)) \geq 0.
\end{align}
An alternative formulation in terms of entropy is:
\begin{align}
    \mathbb{I}(X; Y) = \mathbb{H}(X) - \mathbb{H}(X | Y) 
    = \mathbb{H}(Y) - \mathbb{H}(Y | X).
\end{align}
Mutual information is symmetric and nonnegative. It is zero if and only if $X$ and $Y$ are independent and becomes maximal if and only if $X = Y$. One interprets its value as the amount of information obtained about one random variable while observing another. Since it is defined as an expectation over the joint distribution of the two variables $X$ and $Y$, it is sensitive to higher moments of these distributions, and cannot be captured simply be e.g. their variance.

\subsubsection*{Mutual Information Estimation}
In the context of representation learning, mutual information can be used as a measure of model quality. This is motivated by the data processing inequality, which in this context can be phrased as follows:

Let $X$, $Y$ and $Z$ denote random variables. These are processed as $X \rightarrow Y \rightarrow Z$, in the sense that a sample from $X$ is processed to produce a sample from $Y$, which in turn is processed to produce a sample from $Z$, without any additional knowledge of $X$. Then $Z$ cannot contain more information about $X$ than $Y$, so
\begin{align}
    \mathbb{I}(X;Y) \geq \mathbb{I}(X;Z).
\end{align}
This applied in the context of SBI and NPE, where the parameters of interest $\theta$ are part of the first random variable, the simulated data $x$ part of the second and the summarized data $s_\phi(x)$ part of the third. The data processing inequality then reads:
\begin{align}
    \mathbb{I}(\theta;x) 
    \geq \mathbb{I}(\theta;s_\phi(x)),
\end{align}
and this explains why mutual information is the needed measure of informativeness mentioned in the beginning of this section. A good summary keeps as much information about the original data as possible, subject to the constraint that it should contain less total information than the data itself, i.e. summarize it in some way. The summary distribution is also referred to as the bottleneck.

\red{Something on information bottleneck method?}

In practice, mutual information between the bottleneck $Z$ and the parameter distribution $\Theta$ can only be estimated. In the setting of NPE, one can use the learned posterior to estimate mutual information in a natural way, described in the following. The mutual information between the bottleneck and the parameter distributions is defined as
\begin{align}
    \mathbb{I}(\Theta; Z) &= \int p(\theta, z) \log \frac{p(\theta|z)}{p(\theta)}~\label{eq:MI_int},
\end{align}
where the distribution $p(\theta|z)$ is now estimated by a learned $q_\phi(\theta|z)$, the optimization objective being to minimize the KL-divergence $\mathbb{D}_{\text{KL}}(p||q_\phi)$. The Monte-Carlo approximation of eq.~\ref{eq:MI_int} using this distribution is referred to as the Barber-Agakov estimate~\cite{Barber_2003}:
\begin{align}
    \hat{\mathbb{I}}(\Theta;Z) = \mathbb{I}_{\text{BA}}(\Theta;Z) = \E_{P(\Theta,Z)} \left[ \log q_\phi(\theta|z) - \log p(\theta) \right].
\end{align}
This is computable, because the distribution of the prior $p(z)$ can in practice either be ignored as a constant factor, or in case of SBI is chosen and therefore known. If this is not the case, $p(z)$ can also be estimated separately, which however introduces more uncertainty in the above estimate.

Because the KL-divergence is non-zero, the BA-estimate yields a lower bound on the true mutual information:
\begin{align}
    \mathbb{I}(\Theta;Z) \geq \mathbb{I}_{\text{BA}}(\Theta;Z)
\end{align}

\subsection{An architecture for complex data}

The model discussed in the previous section can be implemented with various different architectural choices. One important choice is that of the \textit{bottleneck size}, that is, the dimensionality of the CNN's output. This part of the model is referred to as the bottleneck because it is the point in which the data are compressed to the highest degree. While a larger bottleneck should be able to capture more information about the data, one needs a principled way to balance the increased model complexity with the increase in information captured by the model. One way to do this is to evaluate the \textit{mutual information} between the feature and the bottleneck space across different such bottleneck dimensions. Mutual information and it's estimation are introduced in the following sections.


An architecture that has proven to be successful in recent SBI approaches \red{some citations here} combines two neural networks in the following principle:
\begin{enumerate}
    \item   The image data are used as input for a convolutional neural network (CNN), which compresses these into a much lower-dimensional summary.
    \item   This summary is mapped to whichever feature is being inferred by a conditional normalizing flow (CNF).
\end{enumerate}
The idea is that the CNN learns to efficiently encode any relevant physics information, while the flow learns to interpret the CNN's output and map it to the inferred quantity. These two networks are introduced in the following.

\subsubsection*{Convolutional Neural Networks}

CNNs are inspired by the task of implementing neural networks in image analysis, i.e. on data with 2D spatial structure. An early version was successfully implemented in 1983 with Fukushima's ”Neucognitron”~\red{Neocognitron citation}. An architecture implementing backpropagation was LeCun's ”LeNet”~\red{LeNet citation}.

To explain the novel concepts behind the CNN in comparison with a MLP, the CNN architecture will first be discussed in 1D. The generalization to higher dimensions will be discussed afterwards.

Let $X$ denote the \red{data space} and $Y$ the feature space. 

Let the input have dimensions $x \in \R^{C\times H \times W}$ where C the channel dimension, H the height dimension and W the width dimension. In case of an RGB image, this channel dimension would be three. 

The CNN learns a function $f_\theta: X \rightarrow Y$. 

As opposed to a simple MLP, the CNN offers the following advantages:

\red{Using summation convention}

\red{Not writing bias explicitly}

A basic MLP consists of only one type of layer. If $h$ is the output of the previous layer, the current layer is defined by
\begin{align}
    z_i = \phi(W_{ij} h_j),
\end{align}
The CNN implements additional layers.

A convolutional layer applies a convolution operaton:

\red{Talk abt convolution vs cross-correlation here}

\red{Is this def even correct? Not need different indices on h?}
\begin{align}
    z_i = \phi(\left[w \circledast h\right]_i) = \phi(\sum_k w_k h_{i+k})
\end{align}
where the dimensionality of the kernel $w\in \R^K$ is typically much lower than that of the data.

A pooling layer applies a pooling operation:
\begin{align}
    z_i = \operatorname{pool} K_{ij}h_j
\end{align}
where $K$ defines a pooling window for each index $i$. Typical choices for $\operatorname{pool}$ are $\operatorname{maxpool}$, which chooses the maximum of all given values, or $\operatorname{meanpool}$, which calculates their mean.

\red{The purpose of a pooling layer is to...}

\red{Add normalization layers here}

\red{Add info on stride params etc here}

\red{Hier kann man schon ein bisschen mehr schreiben, this is all very basic. More info on the CNN specific to this project to come.}

\red{Put it together here, potentially add image of total architecture.}

\red{Also missing MSE loss function as usual loss}

\subsubsection*{Conditional Normalizing Flows}

A normalizing flow is a model that learns to transform a simple (often Gaussian) distribution into a more complex one, the target distribution. The transformation is chosen to be invertible, such that the model becomes generative: once trained, it can be inverted and by transforming samples from the normal, one can sample from the approximated target distribution.

Let $u \sim p_u = \mathcal{N}(0, I_n)$ with $n \in \N$ the base distribution and $x \sim p_x$ with $p_x$ the $n$-dim. target distribution. Let $f: \R^n \to \R^n$ a diffeomorphism, with $g = f^{-1}$ its inverse. The training objective is to find $g$, such that
\begin{align}
    u = g(x) \sim p_u, \qquad x \text{ sampled from } p_x.
\end{align}
The process of generating samples from the approximated distribution is the inverse of the above:
\begin{align}
x = f(u) \sim q_x, \qquad u \text{ sampled from } p_u.
\end{align}
For a given transformation, one computes $p_x(x)$ by applying the change of variables formula:
\begin{align}
q_x(x) = p_u(g(x)) \left| \det \left( \frac{\partial g(x)}{\partial x} \right) \right| = p_u(u) \left| \det \left( \frac{\partial f(u)}{\partial u} \right) \right|^{-1}.
\end{align}
The NLL is the loss function used for optimization, which then becomes:
\begin{align}
    \mathcal{L}_{\text{NF}} = - \E_{x \sim p_x} \Biggl[ \frac{|g(x)|^2}{2} + 
    \log \left| \det \left( \frac{\partial g(x)}{\partial x} \right) \right|\Biggr]
\end{align}
where the first term is a regularization for the latent (in this case Gaussian) space and the second stems from the transformation. See \red{cite Plehns book} for more details.

The normalizing flow becomes expressive by allowing the transformation to be as complex as possible, while restricting it to a class of functions with efficiently computable Jacobian determinants makes it computationally tractable. In practice, one often models $g$ as a composition of simple neural networks, e.g. multilayer perceptrons. \red{How these are invertible...}

\red{Next few lines could also just be skipped}

Let $g_\theta = g_1 \circ g_2 \circ ... \circ g_k$ with $g_i: \R^n \to \R^n, x \mapsto g_i(x | \theta_i)$ and $k \in \N$. 

Now the loss becomes
\begin{align}
    \text{NLL}_\theta =  \sum_{i=1}^{k} \log \left| \det J(g_i) \right| - \log p_u (g_\theta(x))
\end{align}

\red{How to ensure it is a probabilistic model?}

\red{Come up with nicer notation for $g_i (\cdot | \theta_i)$}

\red{Now what this specific Freia flow does. Research on website!}

